\par\medskip
\noindent
13.5. \underline{Caso de una situación relativamente algebraica.}
Suponemos ahora que $A$ procede de un esquema en grupos relativo sobre
$S$, liso y separado, y cuyas fibras son extensiones de esquemas
abelianos por toros. Denotaremos este esquema relativo por la misma
letra $A$. (Nos permitimos este abuso de notación porque el funtor
$G\longmapsto G^{\mathrm{an}}$, de la categoría de los esquemas en
grupos relativos del tipo considerado a la categoría de los grupos
analíticos sobre $S$, es plenamente fiel, como sería fácil demostrar
mediante resultados de tipo GAGA relativo [10].) Para todo lo relativo
a la noción de esquema relativo y a las relaciones entre esquemas
relativos sobre $S$ y espacios analíticos sobre $S$, remitimos a [10]
(véase también [2, Exp. 15]).

Si $\underline O$ denota el anillo local de $S$ en $s$, el esquema
relativo $A$ queda determinado cerca de $s$ por el esquema usual
$A_{\underline O}$ sobre $\operatorname{Spec}(\underline O)$ obtenido
de $A$ por localización en $s$. Es entonces claro que el completado
formal de $A_{\underline O}$ a lo largo de su fibra especial (isomorfa
a $A_s=A_o$) es isomorfo al completado formal del grupo analítico $A$
a lo largo de $A_s$ (13.2), donde, mediante el abuso de notación
introducido arriba, los $A_n=A\times_S S_n$ se identifican con esquemas
algebraicos sobre $S$. En el caso favorable en que $\hat A$, considerado
como esquema formal en grupos sobre $\hat{\underline O}$, es
algebrizable (7.2.1), como sucede en particular cuando $S$ es normal en
$s$ y $A$ es cuasiproyectivo sobre $S$ (7.1.5), se obtiene por tanto un
esquema en grupos
$A^{\natural}_{\hat{\underline O}}$ sobre
$\operatorname{Spec}(\hat{\underline O})$, extensión de un esquema
abeliano $B_{\hat{\underline O}}$ por un toro. Para simplificar,
suponemos que $B_{\hat{\underline O}}$ es proyectivo sobre
$\hat{\underline O}$ (como ocurre, según sabemos, si $\underline O$, y
por tanto $\hat{\underline O}$, es normal). Demostraremos entonces:

\medskip
\noindent
\underline{Proposición 13.5.1.} \underline{En las condiciones
precedentes, el grupo analítico $A^{\natural}$ de (13.3.3) procede, en
un entorno de $s$, de un esquema en grupos relativo determinado salvo
isomorfismo único, que es una extensión de un esquema abeliano relativo
proyectivo sobre $S$ por un toro. Si $A^{\natural}_{\underline O}$
denota el esquema en grupos que este define sobre
$\underline O=\underline O_{S,s}$, entonces el esquema en grupos}

\begin{equation*}
 A^{\natural}_{\underline O}\otimes_{\underline O}
 \hat{\underline O}
\end{equation*}

\noindent
\underline{es canónicamente isomorfo, junto con su estructura de
extensión, a la extensión de Raynaud
$A^{\natural}_{\hat{\underline O}}$ introducida arriba.}

La última afirmación se deduce inmediatamente de la primera, en virtud
de 7.2.1, pues, por (13.3.6), el completado formal de
$A^{\natural}_{\underline O}\otimes_{\underline O}\hat{\underline O}$
a lo largo de su fibra especial es $\hat A_{\underline O}$. Para probar
la algebraicidad de $A^{\natural}$ cerca de $s$, procediendo como en
7.2.1, nos reducimos inmediatamente a probar la del cociente abeliano
$B$ de $A^{\natural}$, lo que se sigue del lema siguiente.

\medskip
\noindent
\underline{Lema 13.5.2.} \underline{Sea $B$ un grupo analítico abeloide
sobre el espacio analítico $S$, y sea $s$ un punto de $S$ tal que el
completado $\hat B$ a lo largo de la fibra $B_s$ sea algebrizable; es
decir, los $B_n=B\times_S S_n$ son algebrizables y $\hat B$,
considerado ahora como esquema formal sobre
$\hat{\underline O}_{S,s}$, es algebrizable, o, equivalentemente,
procede de un esquema abeliano $B_{\hat{\underline O}}$ sobre
$\hat{\underline O}$. Supongamos además que
$B_{\hat{\underline O}}$ es proyectivo sobre $\hat{\underline O}$.
Entonces $B$ es proyectivo sobre $S$ en un entorno de $s$.}

La hipótesis significa que existe una polarización del espacio
analítico formal $\hat B$, es decir, una sucesión de polarizaciones
$p_n$ de los $B_n$ que «se pegan». La conclusión significa que existe
una polarización $p$ de $B$ sobre un entorno de $s$. Sea $P$ el espacio
analítico sobre $S$ de las polarizaciones relativas de $B$ sobre $S$,
que representa el funtor

\begin{equation*}
 S'\longmapsto
 \{\text{polarizaciones de $B_{S'}$ sobre $S'$}\}
\end{equation*}

\noindent
en la categoría de los espacios analíticos $S'$ sobre $S$. Este es un
subfuntor abierto del funtor de Néron--Severi
$\underline{\mathrm{NS}}_{B/S}$ sobre $S$. Los argumentos estándar más
sencillos, que omitimos aquí, muestran que
$\underline{\mathrm{NS}}_{B/S}$, y por tanto también $P$, está
representado por un espacio analítico \underline{no ramificado} sobre
$S$. La hipótesis afirma entonces que existe una sección
\underline{formal} $p$ de $P$ sobre $S$ en $s$. Puesto que $P$ es no
ramificado y, por tanto, localmente cuasifinito sobre $S$, se sigue de
[2, Exp. 18, n.º~1] que toda sección formal de $P$ sobre $S$ en $s$
procede de una sección analítica cerca de $s$, es decir, de una
polarización de $B$ sobre un entorno de $s$. Esto demuestra la
afirmación.

\medskip
\noindent
\underline{Observaciones 13.5.3.}

\smallskip
\noindent
a) Si en 13.5.2 omitimos la hipótesis de que
$B_{\hat{\underline O}}$ sea proyectivo sobre
\sisi{$S$}{$\hat{\underline O}$},
\footnote{La fuente impresa dice aquí «proyectivo sobre $S$»; la
hipótesis del Lema 13.5.2 y el argumento de normalidad que sigue exigen
«proyectivo sobre $\hat{\underline O}$».}
todavía podemos concluir que $B$ es algebrizable cerca de $s$. Para
ello se utiliza un resultado inédito (debido independientemente a
M.~Artin y P.~Deligne, mediante la teoría de los haces $\underline T^i$
de VI 2) que afirma que, si $X$ es propio sobre $S$ y
$X_{\mathrm{red}}$ es algebrizable, entonces $X$ es algebrizable. Esto
nos reduce al caso en que $S$ es reducido. Por otra parte, cuando $S$
es normal, de modo que $\hat{\underline O}$ es normal, se sabe que todo
esquema abeliano sobre $\hat{\underline O}$ es proyectivo sobre
$\hat{\underline O}$ (véase, por ejemplo, [23]), y se aplica 13.5.2.
Para obtener el caso de un $S$ reducido arbitrario, se emplea descenso
no plano desde la normalización $S'$ de $S$, mediante el resultado
siguiente (delicado), que no demostraremos aquí: un morfismo propio de
esquemas noetherianos $f:S'\longrightarrow S$ que sea un epimorfismo
(es decir, tal que
$\underline O_S\longrightarrow f_*(\underline O_{S'})$ sea inyectivo)
es un \underline{morfismo de descenso efectivo} para la categoría
fibrada de los esquemas abelianos relativos.

\smallskip
\noindent
b) Parece que la demostración de 13.5.2 debería probar también el
resultado análogo para todo espacio analítico $B$ propio sobre $S$.
Probablemente sea más delicada la cuestión de si un espacio analítico
$X$ propio sobre $S$, cuyo completado formal $\hat X$ a lo largo de
$X_s$ sea algebrizable, es a su vez algebrizable sobre $S$ cerca de
$s$. Nos limitamos a señalar que aquí la noción de algebraizabilidad
relativa de $X$ sobre $S$ debe tomarse en un sentido más amplio que en
[10]: cerca de $s$, debe exigirse que $X$ proceda de un
\underline{espacio algebraico} sobre $\underline O_{S,s}$ en el sentido
de M.~Artin [1].

\medskip
\noindent
13.5.4. Situémonos en las condiciones de 13.5.1. Para cada entero
$n>\sisi{C}{0}$, elijamos además el isomorfismo estándar de grupos sobre
$\mathbb C$

\begin{equation*}
 \mu_{n\mathbb C}\simeq(\mathbb Z/n\mathbb Z)_{\mathbb C},
 \qquad m\longmapsto\exp(2\pi i m/n)\qquad .
\end{equation*}

\noindent
Al pasar al límite inverso sobre las potencias de un primo fijo
$\ell$, se obtiene un isomorfismo canónico de haces $\ell$-ádicos sobre
$\mathbb C$,

\begin{equation*}
 \mathbb Z_\ell(1)_{\mathbb C}\simeq\mathbb Z_{\ell\mathbb C}
 \qquad ,
\end{equation*}

\noindent
que utilizaremos para identificar ambos miembros.

Aplicando la sucesión exacta de los
$\underline{\operatorname{Ext}}^i_{\mathbb Z}
(\mathbb Z/n\mathbb Z,-)$ a la sucesión exacta (13.1.3), obtenemos un
isomorfismo canónico

\begin{equation}\tag{13.5.5}
 {}_nA^{\mathrm o}\simeq
 \underline R\otimes_{\mathbb Z}\mathbb Z/n\mathbb Z
 \qquad ,
\end{equation}

\noindent
que, cuando $n$ recorre las potencias de un primo dado $\ell$, da un
isomorfismo de sistemas inversos de grupos sobre $S$,

\begin{equation}\tag{13.5.6}
 T_\ell(A^{\mathrm o})\simeq
 \underline R\otimes_{\mathbb Z}\mathbb Z_\ell
 \qquad ,
\end{equation}

\noindent
donde $\mathbb Z_\ell$ se considera como el sistema inverso de los
$\mathbb Z/\ell^{n+1}\mathbb Z$. En la situación aquí considerada, en
la que $A$ procede de un esquema en grupos separado sobre $S$, es
entonces inmediatamente claro que los
$\underline R\otimes\mathbb Z/n\mathbb Z$ pueden considerarse como
esquemas relativos. Esto da sentido al símbolo
$(\underline R\otimes_{\mathbb Z}\mathbb Z_\ell)_{\underline O}$
(¡mientras que $\sisi{R}{\underline R}_{\underline O}$ carecería en general de sentido!), y
proporciona un isomorfismo canónico

\begin{equation}\tag{13.5.7}
 T_\ell(A_{\underline O})\simeq
 (\underline R\otimes_{\mathbb Z}\mathbb Z_\ell)_{\underline O}
 \qquad .
\end{equation}

\noindent
Se sigue inmediatamente de las construcciones del párrafo 2 y de 13.3
que el isomorfismo precedente lleva la parte fija a la parte fija y la
parte tórica a la parte tórica:

\begin{equation}\tag{13.5.8}
 T_\ell(A_{\underline O})^f\simeq
 (\underline R^f\otimes_{\mathbb Z}\mathbb Z_\ell)_{\underline O},
 \qquad
 T_\ell(A_{\underline O})^t\simeq
 (\underline R^t\otimes_{\mathbb Z}\mathbb Z_\ell)_{\underline O}
 \qquad .
\end{equation}

\noindent
Si $A'$ es un grupo analítico relativo sobre $S$ que satisface las
mismas hipótesis que $A$, y si $W_{\underline O}$ es una biextensión de
$(A_{\underline O},A'_{\underline O})$ por
$\mathbb G_{m\underline O}$, y define por tanto sobre un entorno de $s$
una biextensión de $(A,A')$ por $\mathbb G_{mS}$ (utilizando ahora la
noción en sentido analítico complejo, obtenida trabajando en el sitio
de los espacios analíticos sobre $S$ para una topología adecuada),
entonces, por VIII 2.2, $W$ define emparejamientos

\begin{equation}\tag{13.5.9}
 (\underline R\otimes\mathbb Z_\ell)\times
 (\underline R'\otimes_{\mathbb Z}\mathbb Z_\ell)
 \simeq T_\ell(A)\times T_\ell(A')
 \longrightarrow\mathbb Z_\ell(1)_S
 \simeq\mathbb Z_{\ell S}\qquad ,
\end{equation}

\noindent
y es de nuevo claro que el emparejamiento correspondiente

\begin{equation}\tag{13.5.10}
 T_\ell(A_{\underline O})\times T_\ell(A'_{\underline O})
 \longrightarrow\mathbb Z_\ell(1)_S
\end{equation}

\noindent
es precisamente el asociado, por loc. cit., a la biextensión
$W_{\underline O}$ de $(A_{\underline O},A'_{\underline O})$. Tampoco
es difícil ver que, tras restringirse a un entorno de $s$, los
emparejamientos (13.5.9) proceden de un emparejamiento único (que ya no
depende de $\ell$)

\begin{equation}\tag{13.5.11}
 \underline R\otimes\underline R'\longrightarrow\mathbb Z
 \qquad .
\end{equation}

\noindent
Con estos preliminares, los resultados algebraicos relativos a los
$T_\ell(A_{\underline O})$ pueden utilizarse para obtener conclusiones
sobre el haz $\underline R$. Así, el emparejamiento (13.5.11) es cero
sobre $\underline R^t\times\underline R'{}^f$ y sobre
$\underline R^f\times\underline R'{}^t$ (7.4.8). Además, si $S$ es no
singular de dimensión $1$ y, sobre $U=S-\{s\}$, $A$ y $A'$ son esquemas
abelianos relativos duales entre sí mediante $W$, entonces
$\underline R^t_U$ y $\underline R'{}^f_U$, así como
$\underline R^f_U$ y $\underline R'{}^t_U$, son anuladores exactos uno
del otro sobre un entorno suficientemente pequeño de $s$ (7.4.8).
Puesto que (13.5.11) es una dualidad perfecta sobre $U$, se sigue que
hay una dualidad sobre $U$ entre $\underline R/\underline R^t$ y
$\underline R'{}^f$. En particular, $\underline R/\underline R^t$ es
constante y, con mayor razón, la acción de $\pi_1(U,t)$ sobre
$\underline R_t$ es unipotente de nivel $2$. Análogamente, sobre $U$
hay una dualidad perfecta entre $\underline R/\underline R^f$ y
$\underline R'{}^t$. Como este último es también isomorfo al toro dual
de la red $M_S$ sobre $S$, donde $M$ es el grupo definido en (5.5.4),
obtenemos un isomorfismo canónico

\begin{equation}\tag{13.5.12}
 \underline R_U/\underline R^f_U\simeq M_U
 \qquad ,
\end{equation}

\noindent
y el homomorfismo $i$ de (13.3.4) puede escribirse como un homomorfismo

\begin{equation}\tag{13.5.13}
 i:M_U\longrightarrow A^\natural
 \qquad ,
\end{equation}

\noindent
que volveremos a encontrar en un contexto ligeramente distinto en el
párrafo siguiente.
